\documentclass{beamer}
\usepackage[italian]{babel}
\usepackage[utf8]{inputenc}
\usepackage{graphicx}
\usepackage{subfig}
\usepackage{hyperref}
\usepackage{tikz}
\usetikzlibrary{shapes, calc, positioning}
\usepackage{tabularx}
\usepackage{listings}
\usepackage{color}

\definecolor{mygreen}{rgb}{0,0.6,0}
\definecolor{mygray}{rgb}{0.5,0.5,0.5}
\definecolor{mymauve}{rgb}{0.58,0,0.82}

\lstset{ 
  backgroundcolor=\color{black!20!white},   % choose the background color
  basicstyle=\ttfamily,        % the size of the fonts that are used for the code
  captionpos=b,                    % sets the caption-position to bottom
  commentstyle=\color{mygreen},    % comment style
  frame=single,	                   % adds a frame around the code
  %keepspaces=true,                 % useful for keeping indentation 
  language=bash,
  rulecolor=\color{black},         % if not set, the frame-color may be changed on line-breaks within not-black text (e.g. comments (green here))
  showspaces=false,                % show spaces everywhere adding particular underscores; it overrides 'showstringspaces'
  showtabs=false,                  % show tabs within strings adding particular underscores
  tabsize=2,	                   % sets default tabsize to 2 spaces
    deletekeywords={cd},            % if you want to delete keywords from the given language
}

\title{Linux e \texttt{terminali}}
\author{Samuele Musiani, Alice Benatti}
\institute{Università di Bologna, corso di Laurea in Informatica}
\date{14 settembre 2023}
\logo{\includegraphics[width=0.05\textwidth]{assets/by-nc-sa-4-0.png}}

\AtBeginSection[]{
  \begin{frame}
    % \setcounter{tocdepth}{10}
  \vfill
  \centering
  \begin{beamercolorbox}[sep=8pt,center,shadow=true,rounded=true]{title}
    \usebeamerfont{title}\insertsectionhead\par%
  \end{beamercolorbox}
  % \tableofcontents[currentsubsection]
  \vfill
\end{frame}
}

\begin{document}

\begin{frame} 
  \titlepage
\end{frame}

\section{Linux}

\begin{frame}{Perché tutto questo}
  \begin{enumerate}
    \item <1-> Argomenti NON trattati al primo anno
    \item <2-> Estremamente importanti per diventare dei bravi informatici
    \item <3-> Linux è GRATIS
    \item <4-> Linux FUNZIONA BENE
    \item <5-> Linux è open source
  \end{enumerate}
\end{frame}

\begin{frame}{Linux}
  \begin{enumerate}
    \item <1-> Sistema operativo open source sviluppato da Linus Torvalds
    \item <2-> Esistono varie \textit{distribuzioni}: Debian, Fedora, Arch, ecc.
    \item <3-> Include una \textit{shell} potentissima
    \item <4-> Rispetto a Windows ha una gestione delle risorse estremamente più
      efficiente
    \item <5-> Estremamente affidabile e usato in tutti gli ambiti server
    \item <6-> Le varie distribuzioni hanno ottimi package manager!
  \end{enumerate}
\end{frame}

\begin{frame}{Intallare per la prima volta Linux}
  Esistono vari modi per utilizzare e provare Linux. I due principali sono:
  \begin{enumerate}
    \item <1-> Installarlo in \textbf{dual-boot}: Se avete un pc con installato
      Windows potete installarci "a fianco" una distribuzione Linux. 

      \only<1> {\bigskip ATTENZIONE: non sto parlando della WSL (Windows 
      Subsystem Linux)! \textbf{NON USATE LA WSL}.}

    \item <2-> Installare Linux in una \textbf{virtual machine} tramite 
      VirtualBox. È come avere un pc virtuale su cui provare Linux. Avete 
      il vostro sistema principale che consuma risorse. Unica soluzione per chi 
      usa macOS.
  \end{enumerate}
\end{frame}

\begin{frame}{Quale distribuzione scegliere?}
  Se non sapete cosa installare per la prima volta vi consiglio la distribuzione
  \textbf{Debian} con l'ambiente grafico \textbf{Gnome}.\bigskip

  La distribuzione determina principalmente il \textit{package manager} che 
  usate. Debian usa \texttt{apt}.\bigskip

  Non installate Ubuntu! Nonostante sia basato su Debian non è sicuramente una
  delle scelte migliori.\bigskip

  I use Arch btw.
\end{frame}

\section{Shell e terminali}

\subsection{Cos'è una shell}
\begin{frame}[fragile]{Cos'è una shell}
  Una \texttt{shell} è un programma che permette di parlare con il sistema 
  operativo attraverso dei comandi da tastiera.
  \pause

  \begin{figure}
        \begin{lstlisting}
[samu@leibniz ~]$
    \end{lstlisting}
    \caption{Esempio di un \textit{shell prompt}}
  \end{figure}
\end{frame}

\subsection{Perché tornare all'età della pietra}
\begin{frame}{Perché tornare all'età della pietra}
  I primi computer utilizzavano i terminali, perché dovremmo farlo anche noi
  quando abbiamo un'ottima interfaccia grafica?\bigskip
  \pause

  \textit{"Graphical user interface make easy taks easy, while command line 
  iterfaces make difficult tasks possible"}
\end{frame}

\subsection{Vantaggi}
\begin{frame}{Vantaggi di un terminale}
  I primi computer utilizzavano i terminali, perché dovremmo farlo anche noi
  quando abbiamo un'ottima interfaccia grafica?\bigskip

  \begin{itemize}
    \item <1-> È decisamente più veloce di una GUI
    \item <2-> Avete il completo controllo di quello che state facendo
    \item <3-> Moltissime cose non si possono fare con una GUI
    \item <4-> Il vero potere della \texttt{shell} sono le \textit{pipe}
    \item <5-> Accedere a server e configurare servizi
  \end{itemize}
\end{frame}

\section{I primi passi in un terminale}
\begin{frame}{Aprire un terminale}
  Per utilizzare una \texttt{shell} è necessario disporre di un emulatore di 
  terminale.\bigskip

  Un emulatore di terminale è un'applicazione che permette di interagire con
  la \texttt{shell}.
\end{frame}

\subsection{Digitare e leggere}
\begin{frame}[fragile]{Digitare e leggere}
  Proviamo a digitare delle lettere a caso nella nostra \texttt{shell}:
  \begin{figure}
    \begin{lstlisting}
[samu@leibniz ~]$ fasdjfivb
bash: fasdjfivb: command not found
[samu@leibniz ~]$ 
    \end{lstlisting}
    \caption{Esempio di lettere digitate a caso in una \textit{shell}}
  \end{figure}
  \pause
  La \texttt{shell} non ha riconosciuto il comando e ce lo ha scritto a schermo.
\end{frame}

\begin{frame}[fragile]{Digitare e leggere}
  Proviamo ora a digitare un comando esistente:
  \begin{figure}
        \begin{lstlisting}
[samu@leibniz ~]$ date
Thu Set 7 12:15:58 PM CEST 2023
[samu@leibniz ~]$
    \end{lstlisting}
    \caption{Esempio di lettere digitate a caso in una \textit{shell}}
  \end{figure}
  \pause
  In questo caso la \texttt{shell} ha riconosciuto il comando e chi ha stampato
  la data, l'ora e qualche altra informazione come il fuso orario.
\end{frame}

\subsection{Quali comandi esistono}
\begin{frame}{Quali comandi esistono}
  Come abbiamo appena visto alcuni comandi esistono e altri no.\bigskip

  Non vi è una lista dei comandi esistenti perché sarebbe troppo lunga.\bigskip

  Oggi vi mostreremo alcuni comandi di base per iniziare a familiarizzare con la
  \texttt{shell}.
\end{frame}

\section{Muoversi tra i file e le cartelle}

\subsection{Elencare i file - ls}
\begin{frame}[fragile]{Elencare i file}
  Le operazioni più importanti sono legate alla gestione dei file. Proviamo a 
  digitare il comando \texttt{ls} nella shell.\bigskip

    \begin{lstlisting}
[samu@leibniz ~]$ ls
Desktop Documents Downloads Music Pictures 
Public  Templates Video
[samu@leibniz ~]$
    \end{lstlisting}

  Apriamo ora il programma per gestire i file a livello grafico e vediamo che
  solo gli stessi.
\end{frame}

\subsection{Muoversi tra le cartelle - cd}
\begin{frame}[fragile]{Entrare in una cartella}
  Dall'interfaccia grafica possiamo entrare nelle cartelle. Con la \texttt{shell}
  possiamo farlo attraverso il comando \texttt{cd}.\bigskip

  Per farlo però è necessario indicare in quale cartella ci vogliamo muovere 
  visto che potrebbe essercene più di una.\bigskip

  Per farlo è quindi necessario fornire al comando \texttt{cd} un 
  \textbf{argomento}.\bigskip

    \begin{lstlisting}
[samu@leibniz ~]$ cd Documents
[samu@leibniz ~/Documents]$ cd Documents
    \end{lstlisting}
\end{frame}

\begin{frame}[fragile]{Entrare in una cartella}
    \begin{lstlisting}
[samu@leibniz ~]$ cd Documents
[samu@leibniz ~/Documents]$ cd Documents
    \end{lstlisting}

  L'argomento è successivo al comando. Un modo per pensare a questo ordine è
  chiedersi:\medskip
  
  \begin{itemize}
    \item "Cosa voglio fare?" -> Entrare di una cartella -> \texttt{cd}
    \item "In quale cartella voglio entrare?" -> \texttt{cd} 
      \textit{nomecartella}
  \end{itemize}
\end{frame}

\begin{frame}[fragile]{Uscire da una cartella}
  Per uscire da una cartella il comando è \texttt{cd ..}\bigskip

    \begin{lstlisting}
[samu@leibniz ~/Documents]$ cd ..
[samu@leibniz ~]$
    \end{lstlisting}

  L'argomento \texttt{..} indica sempre la cartella precedente a quella attuale.
\end{frame}

\begin{frame}{Current working directory}
  In Linux le cartelle si chiamano \textit{directory}. Esiste un comando per
  stampare il \textit{path} della cartella corrente: \texttt{pwd}.\bigskip

  Il path assume la seguente forma: \texttt{/home/samu}.\bigskip

  Il carattere \texttt{/} viene utilizzato come separatore.
\end{frame}

\begin{frame}{Current working directory}
  Il path assume la seguente forma: \texttt{/home/samu}.\bigskip

  Notiamo che il \textit{path} inizia con uno \texttt{/}\bigskip

  In Linux la radice del file-system è proprio \texttt{/}
\end{frame}

\section{Linux file system}

\subsection{Dove tutto ha inizio}
\begin{frame}{Dove tutto ha inizio}
  Per fare riferimento a uno file è necessario identificarlo con un 
  \textit{path}. \bigskip

  I \textit{path} sono unici, non possono esserci quindi file diversi con lo 
  stesso \textit{path}. \bigskip

  Ne consegue che se due file hanno lo stesso \textit{path} sono lo stesso file.
  \bigskip

  Ogni \textit{path} in linux inizia sempre dalla radice del file-system, ovvero
  \textit{/}\bigskip

  Tutte le cartelle presenti in \texttt{/} sono riservate al sistema e al suo
  corretto funzionamento
\end{frame}

\begin{frame}[fragile]{Cartelle di sistema}
  \begin{figure}
    \begin{lstlisting}[basicstyle=\footnotesize]
[samu@leibniz /]$ ls
bin   dev  home  lib64       mnt  proc  run   tmp  var
boot  etc  lib   lost+found  opt  root  sbin  sys  usr
[samu@leibniz /]$
    \end{lstlisting}
    \caption{Lista dei file presenti in \texttt{/}}
  \end{figure}
  \begin{itemize}
    \item \texttt{/bin}: contiene programmi necessari al sistema per funzionare
    \item \texttt{/boot}: contiene il kernel e altri file necessari al sistema
      per partire.
    \item \texttt{/etc}: contiene tutti i file di configurazione del sistema.
    \item \texttt{/home}: contiene le cartelle riservate agli utenti
    \item \texttt{/tmp}: contiene file temporanei che vengono cancellati ad ogni
      spegnimento del sistema
    \item \texttt{/usr}: contiene programmi e file usati dagli utenti
  \end{itemize}
\end{frame}

\subsection{\textit{Path} relativi e assoluti}
\begin{frame}{\textit{Path} relativi e assoluti}
  Per identificare un file è possibile usar due tipi di \textit{path}: assoluto 
  e relativo:
  \begin{enumerate}
    \item Un \textit{path} assoluto è un percorso ad un file che inizia da 
      \texttt{/} e termina con il nome di quel file. prende quindi le seguenti
      forme:
      \begin{itemize}
        \item \texttt{/home/samu/slides.tex}
        \item \texttt{/usr/bin/firefox}
        \item \texttt{/tmp}
      \end{itemize}
    \item Un \textit{path} relativo è il percorso necessario per raggiungere un
      file rispetto alla \textbf{cartella corrente}. prende quindi le seguenti
      forme:
      \begin{itemize}
        \item \texttt{slides.tex}
        \item \texttt{../}
        \item \texttt{immagini/gattini.png}
      \end{itemize}
  \end{enumerate}
\end{frame}

\section{Interagire con file}

  \subsection{Distinguere i tipi di file - \texttt{file}}
\begin{frame}[fragile]{A chi serve nautilus?}
  Per sapere di che tipo è un file possiamo usare il comando \texttt{file}:
  \begin{figure}
    \begin{lstlisting}[basicstyle=\footnotesize]
[samu@leibniz ~]$ file Documents/
Documents/: directory
[samu@leibniz ~]$ file screen.png
screen.png: PNG image data, 3840x1080, 8-bit/color RGB
[samu@leibniz ~]$ file lista
lista: Unicode text, UTF-8 text
[samu@leibniz ~]$ file /bin/firefox
/bin/firefox: POSIX shell script, ASCII text executable
[samu@leibniz ~]$
    \end{lstlisting}
    \caption{esempi del comando file}
  \end{figure}
  In linux i file non hanno bisogno di un'estensione, è quindi molto utile 
  questo comando per determinare il tipo di un file
\end{frame}

\subsection{Copiare, spostare ed eliminare}
\begin{frame}[fragile]{Copiare file}
  Il comando per copiare dei file è \texttt{cp}. il suo utilizzo è 
  principalmente: \texttt{cp file/da/copiare destinazione}
  \begin{figure}
    \begin{lstlisting}[basicstyle=\small]
[samu@leibniz prova]$ ls
castoro  ligma  piante
[samu@leibniz prova]$ cp castoro criceto
[samu@leibniz prova]$ ls
castoro  criceto  ligma  piante
[samu@leibniz ~]$
    \end{lstlisting}
    \caption{esempi del comando cp}
  \end{figure}
  Sono ammessi sia \textit{path} assoluti sia relativi
\end{frame}

\begin{frame}[fragile]{Spostare file}
  Il comando per spostare (tagliare) dei file è \texttt{mv}. il suo utilizzo è 
  pressoché identico al comando di copia: \texttt{mv file/da/muovere 
  destinazione}
  \begin{figure}
    \begin{lstlisting}[basicstyle=\small]
[samu@leibniz prova]$ ls
castoro  criceto  ligma  piante
[samu@leibniz prova]$ cd piante/
[samu@leibniz prova]$ ls
cubo
[samu@leibniz piante]$ mv ../castoro .
[samu@leibniz piante]$ ls
castoro  cubo
    \end{lstlisting}
    \caption{esempi del comando mv}
  \end{figure}
  La destinazione \texttt{.} indica la directory corrente.
\end{frame}

\begin{frame}[fragile]{Rinominare file}
  Il comando \texttt{mv} permette anche di rinominare i file:
  \begin{figure}
    \begin{lstlisting}[basicstyle=\normalsize]
[samu@leibniz piante]$ ls
castoro  cubo
[samu@leibniz piante]$ mv cubo triangolo
[samu@leibniz piante]$ ls
castoro  triangolo
    \end{lstlisting}
    \caption{esempi del comando mv}
  \end{figure}
\end{frame}

\begin{frame}{Now I am become death, the destroyer of files}
  Il comando più pericoloso in linux è indubbiamente \texttt{rm}.\bigskip

  \texttt{rm nome/file} elimina il file passato come argomento.\bigskip

  Non si può tornare indietro, una volta eliminato un file è perso per 
  sempre!\bigskip

  La \textit{shell} si aspetta che voi sappiate esattamente quello che state 
  facendo e non si preoccupa se questo può distruggere il sistema.
\end{frame}

\begin{frame}[fragile]{Copiare una cartella}
  Per copiare le directory il comando \texttt{cp} deve funzionare in modalità 
  \textit{ricorsiva}, per permettere la copia di tutti gli elementi all'intero
  della directory.\bigskip

  Per copiare una cartella è quindi necessario aggiungere la \textit{flag} 
  \texttt{-r}
  \begin{figure}
    \begin{lstlisting}[basicstyle=\footnotesize]
[samu@leibniz prova]$ ls
criceto  ligma  piante
[samu@leibniz prova]$ cp piante roveri
cp: -r not not specified; omitting directory 'piante'
[samu@leibniz prova]$ cp -r piante roveri
[samu@leibniz prova]$ ls
criceto  ligma  piante  roveri
    \end{lstlisting}
    \caption{copia di una cartella}
  \end{figure}
\end{frame}

\subsection{Flags - man}
\begin{frame}{Flags}
  Le \textit{flag} sono un modo per estendere le funzionalità di un 
  comando.\bigskip

  Vengono specificate dopo il comando e sono precedute da un trattino 
  \texttt{-}\bigskip

  Per comodità sono di una sola lettera, ma in certi casi possono essere anche 
  più verbose es. \texttt{--recursive}.\bigskip

  Si possono combinare più \textit{flag} concatenando le lettere dopo il 
  trattino: \texttt{-r -t} è equivalente a \texttt{-rt}
\end{frame}

\begin{frame}{Man}
  Esiste un comando per leggere il \textit{manuale} di un comando: 
  \texttt{man}.\bigskip

  Sintassi: \texttt{man comando}\bigskip

  Aprirà un lettore di testo integrato nel terminale chiamato \texttt{less}. Per 
  navigare sono usati i seguenti comandi:
  \begin{itemize}
    \item \texttt{j}: Muoversi verso il basso
    \item \texttt{k}: Muoversi verso l'alto
    \item \texttt{g}: Inizio del file
    \item \texttt{G}: Fine del file
    \item \texttt{/name}: Cerca la stringa \texttt{name} in tutto il testo
    \item \texttt{n}: Selezionata una stringa va all'occorrenza successiva
    \item \texttt{q}: Esce e torna al promp
  \end{itemize}
\end{frame}

\subsection{Lettura di file - less}
\begin{frame}{Less}
  Esiste un lettore di testo integrato nella \textit{shell} chiamato 
  \texttt{less}.\bigskip

  \texttt{less} è in realtà un comando che permette di leggere file di 
  testo.\bigskip

  Sintassi: \texttt{less nome/file}\bigskip

  \texttt{less} è molto veloce a leggere file di testo di grandi dimensioni.
\end{frame}

\section{Operazioni testuali}
\begin{frame}{Contare caratteri, linee, ecc. - wc}
  Esiste un comando per contare i caratteri, le linee e altre informazioni
  all'interno di un file di testo: \texttt{wc}\bigskip

  Sintassi: \texttt{wc nome/file}\bigskip

  Senza nessuna \textit{flag} stampa:
  \begin{itemize}
    \item Il numero di righe
    \item Il numero di parole
    \item Il numero di bytes
  \end{itemize}
  Per stampare il numero di caratteri: \texttt{wc -c nome/file}
\end{frame}

\begin{frame}[fragile]{Ricerca di una stringa - grep}
  Per cercare una stringa in un file esiste \texttt{grep}\bigskip

  Sintassi: \texttt{grep "string" file}\bigskip

  Esiste la \textit{flag} \texttt{-i} per la ricerca case-insensitive.\bigskip

  Uno dei comandi più potenti per la ricerca di pattern.
  \begin{figure}
    \begin{lstlisting}[basicstyle=\footnotesize]
[samu@leibniz castorini]$ grep "13" LICENSE
 13. Use with the GNU Affreo General Pubic License.
Section 13, concerning interaction through a network
[samu@leibniz castorini]$
    \end{lstlisting}
    \caption{Ricerca di una stringa}
  \end{figure}
\end{frame}

\begin{frame}{Concatenazione o lettura? - cat}
  Il comando \texttt{cat} è nato per concatenare più file.\bigskip

  Per renderlo però completamente funzionante abbiamo bisogno dell'operatore
  di ridirezione che verrà spiegato più avanti.\bigskip

  Possiamo però usarlo per leggere file generalmente corti. Spesso è più rapido
  da usare di \texttt{less}.\bigskip

  A differenza di \texttt{less} stampa il file completo sul terminale.
\end{frame}

\begin{frame}{Testa e coda - head, tail}
  Il comando \texttt{cat} prende uno o più file in input e li stampa sul 
  terminale.\bigskip

  Esiste anche \texttt{head} che stampa solo le prime 10 righe di un file (si
  possono modificare con la \textit{flag} \texttt{-n numero}).\medskip

  Da notare che la \textit{flag} precedente ha preso un parametro.\bigskip

  Esiste anche \texttt{tail} che fa esattamente la stessa cosa di \texttt{head}, 
  ma partendo dalla fine del file.
\end{frame}

\section{Complementi di comandi base}

\begin{frame}{Complementi di comandi base}
  Per pulire il terminale esiste il comando \texttt{clear}\bigskip

  Per resettare il terminale esiste il comando \texttt{reset}. Sarà molto utile
  quando lavorerete al progetto di programmazione e romperete tutto con 
  la libreria grafica.\bigskip

  Per vedere i vecchi comandi eseguiti esite il comando \texttt{history}\bigskip

  Per riprendere un comando eseguito di recente basta utilizzare la freccetta
  verso l'alto.\bigskip

  Per editare un comando si possono usare le freccette verso destra e verso 
  sinistra.\bigskip
\end{frame}

\begin{frame}{Copia, incolla e interruzione}
  Nel terminale non funziona il classico copia e incolla da tastiera eseguito
  con \texttt{ctrl + c} e \texttt{ctrl + v}. Queste combinazioni di tasti hanno 
  il loro scopo e non sono fatti per copiare.\bigskip

  Per copiare e incollare dovete usare \texttt{shift + ctrl + c} e 
  \texttt{shift + ctrl + v}.\bigskip

  \texttt{ctrl + c} serve per interrompere un processo in esecuzione.
\end{frame}

\section{Pipe e ridirezione: il potere della shell}
\begin{frame}{Il potere della shell}
  Finora abbiamo visto soltanto comandi base e usati singolarmente.\bigskip

  Nonostante molti di essi siano utili anche da soli, sono sicuramente più 
  utili usati in combinazione con altri comandi.\bigskip

  Come si fa però a combinare più comandi? \pause
  Ricordiamoci che:
  \begin{itemize}
    \item <2-> Un comando restituisce sempre qualcosa sullo standard output
    \item <3-> Lo standard output è considerato come un file dal sistema
    \item <4-> La maggior parte dei comandi visti fino ad adesso hanno la 
      possibilità di prendere in input lo standard output invece che un file 
      classico
  \end{itemize}
\end{frame}

\begin{frame}{Il potere della shell}
  \begin{figure}
    \begin{tikzpicture}
      \draw (-1.5,3) rectangle ++(3, 1);
      \draw (-1.5,-3) rectangle ++(3, 1);
      \draw [->] (0,3) -- (0, 1);
      \draw [->] (0,-0.9) -- (0, -2);
      \node [cylinder, shape border rotate=90, draw,minimum height=2cm,minimum width=1.5cm] {};

      \node {buffer};
      \node at (0,3.5) {command 1};
      \node at (0,-2.5) {command 2};
      \node at (1, 2) {output};
      \node at (1,-1.5) {input};
    \end{tikzpicture}
  \end{figure}
\end{frame}

\begin{frame}{Pipe}
  Per prendere l'output di un comando e riderizionarlo in input verso un altro
  comando si usa la \textbf{pipe} \texttt{|}\bigskip

  Per esempio se vogliamo vedere tutti i file presenti in \texttt{/bin} il 
  nostro terminale si riempie di scritte.\bigskip

  Possiamo visualizzare il lungo output con il comando \texttt{less}: 
  \texttt{ls /bin | less}
\end{frame}

\begin{frame}{Pipe - grep}
  Come abbiamo visto la lista di file presenti in \texttt{/bin} è molto lunga.
  Se volessimo trovarne uno specifico?\pause \bigskip

  Nonostante esiste un comando apposito per cercare file, possiamo fare:
  \texttt{ls /bin | grep "firefox"} dove al posto di \textit{firefox} può 
  andarci una qualsiasi stringa.
\end{frame}

\begin{frame}{Pipe - esempi}
  Di seguito una lista di esempi di utilizzo della \texttt{pipe}:
  \begin{itemize}
    \item \texttt{cat file1 file2 | grep "word"} cerca una stringa in più file
    \item \texttt{ls /bin | wc -l} conta quanti programmi sono presenti in 
      \texttt{/bin}
    \item \texttt{ls /bin | grep "zip" | wc -l} conta quanti programmi hanno
      la stringa "\textit{zip}" al loro interno nella cartella \texttt{/bin}
    \item \texttt{grep "castoro" animali | wc -l} conta le occorrenze di 
      \texttt{castoro} trovate nel file \texttt{animali}
    \item \texttt{grep "the" book | less} mostra le occorrenze di \texttt{the} 
      trovate in \texttt{book} attraverso il lettore \texttt{less}
  \end{itemize}
\end{frame}

\begin{frame}{Ridirezione su file}
  Come abbiamo visto è possibile mandare l'output di un comando nell'input di
  un altro comando.\bigskip

  Se volessimo salvare l'output di un comando su un file?\pause

  Esiste l'\textbf{operatore di ridirezione} \texttt{>}\bigskip

  Al posto di indirizzare l'output in un comando, scrive direttamente su un
  file.\bigskip

  Sintassi: \texttt{comando > file}\bigskip

  ATTENZIONE: Alla \textit{shell} non interessa se il file esiste già, quindi se 
  esiste lo SOVRASCRIVE COMPLETAMENTE.
\end{frame}

\begin{frame}{Ridirezione su file non distruttiva}
  Esiste anche un operatore per indirizzare su file l'output di un comando senza 
  sovrascrivere il contenuto del file, ma "appendendo" alla fine del file il
  contenuto scritto.\bigskip

  Sintassi: \texttt{comando >> file}\bigskip

  Si usa nello stesso modo dell'operatore classico\bigskip
\end{frame}


\section{Wildcard}
\begin{frame}{Wildcard}
  Nella \textit{shell} l'asterisco \texttt{*} fa da "segnaposto" per una 
  qualsiasi altra sequenza di caratteri.\bigskip

  Esempio: \texttt{torr*} si espande in: \texttt{torr}, \texttt{torra}, 
  \texttt{torrb}, $\cdots$, \texttt{torraa}, \texttt{torrab}, \texttt{torrac}, 
  $\cdots$\bigskip

  Valgono ovviamente anche i numeri e altri caratteri oltre alle lettere.\bigskip

  Questo permette di indicare più file con parti comuni nel nome.\bigskip

  Si possono combinare anche più asterischi: \texttt{c*a*} fa match con tutte le
  parole che iniziano per \texttt{c} e hanno almeno una \texttt{a} nel 
  nome.\bigskip

  \texttt{*pila*} fa match con tutte le parole che hanno \texttt{pila} nel 
  nome.\bigskip

\end{frame}

\section{Root e permessi}

\subsection{Sistema multi-utente}
\begin{frame}{Sistema multi-utente}
  Linux è un sistema \textit{multi-utente}. Più persone possono usare lo stesso
  computer simultaneamente.\bigskip

  Questo utilizzo simultaneo non è da intendere con più schermi, tastiere e mouse
  attaccati allo stesso dispositivo fisico.\bigskip

  Si può utilizzare un computer connesso alla rete tramite il comando 
  \texttt{ssh}.\bigskip

  \texttt{ssh} stabilisce una connessione cifrata tra l'utente e il computer e
  fornisce una \textit{shell} su cui poter lavorare.\bigskip

  È quindi importante familiarizzare con la \textit{shell} in quanto è l'unico
  modo per amministrare macchine e server remoti.
\end{frame}

\begin{frame}{Sistema di permessi}
  La possibilità per Linux di gestire più utenti non è un aggiunta recedente. Il 
  sistema è stato scritto proprio per supportare questa funzionalità.\bigskip

  Essendoci quindi più utenti è necessario avere un sistema di permessi 
  adeguato.\bigskip

  Il sistema di permessi utilizzato da Linux non è semplice, e di seguito sarà
  data soltanto un'introduzione.\bigskip
\end{frame}

\begin{frame}{Root}
  In tutti i sistemi Linux esiste un unico utente che ha i permessi per seguire
  qualsiasi operazione: \textbf{root}\bigskip

  È l'amministratore del sistema.\bigskip

  NON deve essere MAI usato come utente se non per le operazioni strettamente
  necessarie.\bigskip

  Per aprire una \textit{shell} come utente \texttt{root} è possibile digitare
  il comando: \texttt{su}\bigskip

  Il comando chiederà quindi la password di \texttt{root} (che generalmente è
  impostata durante l'installazione del sistema) e se corretta aprirà una 
  \textit{shell} con i privilegi di amministratore.
\end{frame}

\begin{frame}{Permessi su un file}
  Dentro la cartella \texttt{/etc} esiste un file chiamato \texttt{shadow}. 
  Se proviamo a leggerlo con \texttt{less} otteniamo: \textbf{/etc/shadow: 
  Permission denied}\bigskip

  Questo significa che il nostro utente non ha i permessi per leggere il file.
  Ma come potevamo saperlo a priori senza tentare di leggerlo?\pause\bigskip

  Una \textit{flag} molto usata per il comando \texttt{ls} è 
  \texttt{-l}.\bigskip

  \texttt{ls -l} permette di vedere molte più informazioni sui file presenti in
  una directory.
\end{frame}

\begin{frame}{Esempio di ls -al}
  \begin{figure}
    \begin{tikzpicture}
      \node[inner xsep=0pt,brown] (a) {\texttt{d}};

      \node[inner xsep=0pt,right=0pt of a, blue] (b) {\texttt{rwxr-xr-x}};
      \node[inner xsep=0pt,right=5pt of b, green!50!black] (c) {\texttt{samu}};
      \node[inner xsep=0pt,right=5pt of c, orange] (d) {\texttt{samu}};
      \node[inner xsep=0pt,right=5pt of d] (e) 
      {\texttt{4096 Sep  6 15:02 Documents}};

      \draw [brown] (a.south west) -- (a.south east);
      \draw [blue] (b.north west) -- (b.north east);
      \draw [green!50!black] (c.south west) -- (c.south east);
      \draw [orange] (d.north west) -- (d.north east);

      \coordinate (a1) at (a.south);
      \coordinate (b1) at (b.north);
      \coordinate (c1) at (c.south);
      \coordinate (d1) at (d.north);

      \draw[brown] ($(a1)+(-90:0.1)$)--($(a1)+(-90:1)$) 
      node[fill=white,font=\normalsize\sffamily,inner sep=2pt] 
      {Directory o file};

      \draw[blue] ($(b1)+(90:0.1)$)--($(b1)+(90:1.5)$) 
      node[fill=white,font=\normalsize\sffamily,inner sep=5pt, text 
      width=3.5cm] {Permessi di lettura, scrittura ed esecuzione};

      \draw[green!50!black] ($(c1)+(-90:0.1)$)--($(c1)+(-60:2)$) 
      node[fill=white,font=\normalsize\sffamily,inner sep=2pt] 
      {User owner};

      \draw[orange] ($(d1)+(90:0.1)$)--($(d1)+(50:2)$) 
      node[fill=white,font=\normalsize\sffamily,inner sep=5pt,] 
      {Group owner};
    \end{tikzpicture}
  \end{figure}
\end{frame}

\begin{frame}{Read, write, execute}
  Ogni file ha un stringa formata da \textbf{9 bit} che determina quali 
  permessi specifici ha quel file rispetto all'\textit{utente}, il 
  \textit{gruppo} e gli \textit{altri}.\bigskip

  I 9 bit sono suddivisi in \textbf{gruppi di 3}: il primo è specifico per 
  l'\textit{utente}, il secondo è specifico per il \textit{gruppo} e i rimanenti
  sono per tutti gli \textit{altri}.
  \begin{figure}
    \begin{tikzpicture}
      \node[inner xsep=0pt,orange] (a) {\texttt{rwx}};

      \node[inner xsep=0pt,right=10pt of a, blue] (b) {\texttt{rwx}};
      \node[inner xsep=0pt,right=10pt of b, green!50!black] (c) {\texttt{rwx}};

      \draw [brown] (a.south west) -- (a.south east);
      \draw [blue] (b.north west) -- (b.north east);
      \draw [green!50!black] (c.south west) -- (c.south east);

      \coordinate (a1) at (a.south);
      \coordinate (b1) at (b.north);
      \coordinate (c1) at (c.south);

      \draw[brown] ($(a1)+(-90:0.1)$)--($(a1)+(-90:1)$) 
      node[fill=white,font=\normalsize\sffamily,inner sep=2pt] 
      {User};

      \draw[blue] ($(b1)+(90:0.1)$)--($(b1)+(90:1)$) 
      node[fill=white,font=\normalsize\sffamily,inner sep=5pt] {Group};

      \draw[green!50!black] ($(c1)+(-90:0.1)$)--($(c1)+(-90:1)$) 
      node[fill=white,font=\normalsize\sffamily,inner sep=2pt] 
      {Other};
    \end{tikzpicture}
  \end{figure}
\end{frame}

\begin{frame}{Tabella con i permessi}
  \begin{table}
    \begin{tabularx}{\textwidth}{|l|l|X|}
      \hline
      & File  & Directory  \\ \hline
      r & Permette la lettura di un file & Permette di vedere il contenuto se 
      anche \texttt{x} è segnato\\ \hline
      w & Permette di scrivere sul file  & Permette di creare ed eliminare file 
      dentro la directory se \texttt{x} è segnato \\ \hline
      x & Permette di eseguire un file   & Permette di entrare nella 
      directory \\ \hline
    \end{tabularx}
    \caption{Significato dei permessi per file e directory}
  \end{table}
\end{frame}

\begin{frame}{Cambiare owner e group}
  Per cambiare l'owner di un file e il gruppo possiamo usare il comando 
  \texttt{chown}.\bigskip

  Sintassi: \texttt{chown user:group file}\bigskip

  Per eseguire il comando sono necessari i privilegi di root.
\end{frame}

\begin{frame}{Cambiare permessi wrx}
  Per cambiare i permessi lettura, scrittura ed esecuzione si utilizza il 
  comando \texttt{chmod}.\bigskip

  Sintassi: \texttt{chmod [PART][ACTION][PERMISSION] file}\bigskip

  \only<1> {
    Al posto di \texttt{[PART]} è necessario specificare la 
    parte che si vuole modificare:
    \begin{itemize}
      \item user: si utilizza \texttt{u}
      \item group: si utilizza \texttt{g}
      \item others: si utilizza \texttt{o}
      \item all: si utilizza \texttt{a}
    \end{itemize}
  }

  \only<2> {
    Al posto di \texttt{[ACTION]} è necessario specificare la 
    l'azione da compiere
    \begin{itemize}
      \item \texttt{+}: Aggiunge il permesso
      \item \texttt{-}: Rimuove il permesso
      \item \texttt{=}: Assegna esattamente quel permesso
    \end{itemize}
  }

  \only<3> {
    Al posto di \texttt{[PERMISSION]} è necessario specificare la 
    il permesso o i permessi da moficare:
    \begin{itemize}
      \item \texttt{r}: Read
      \item \texttt{w}: Write
      \item \texttt{x}: Execute
    \end{itemize}
  }

  \bigskip Per eseguire il comando sono necessari i privilegi di root.
\end{frame}

\begin{frame}{Esempi \texttt{chmod}}
  \begin{itemize}
    \item \texttt{chmod u+x pippo} Rende pippo un file eseguibile per l'utente
    \item \texttt{chmod o-w pippo} Rimuove la possibilità a tutti gli utenti
      diversi dall'owner del file e non appartenenti al gruppo del file di 
      scrivere su pippo.
    \item \texttt{chmod g+r pippo} Rende pippo leggibile al gruppo
    \item \texttt{chmod g+x pippo} Rende pippo eseguibile dal gruppo
    \item \texttt{chmod u=rwx,g=,o= pippo} Rende pippo leggibile, scrivibile ed
      eseguibile per l'utente. Inoltre rimuove tutti i permessi dal gruppo e 
      altri.
  \end{itemize}
\end{frame}

\section{More work to do}
\begin{frame}{More work to do}
  I comandi presentati sono soltanto una piccolissima parte dell'infinità di 
  comandi presenti in una sistema Linux.\bigskip 

  Non possiamo ovviamente includerli tutti, ma di seguito lasceremo alcuni 
  comandi che potete approfondire:
  \begin{itemize}
    \item \texttt{ssh}
    \item \texttt{sudo}
    \item \texttt{apt}, \texttt{yum} e \texttt{pacman} (Dipende dalla 
      distribuzione Linux)
    \item \texttt{top} e \texttt{htop}
    \item \texttt{kill}
    \item \texttt{touch}, \texttt{locate} e  \texttt{find} 
    \item \texttt{nano} e \texttt{vim}
  \end{itemize}
\end{frame}
    
\end{document}
